
% Thesis Abstract -----------------------------------------------------


%\begin{abstractslong}    %uncommenting this line, gives a different abstract heading
\begin{abstracts}        %this creates the heading for the abstract page




{\color{red} 
Value-based reinforcement learning, pioneered by Deep Q-Networks, established that agents could learn complex behaviours directly from raw sensory input. Actor-critic methods, developed subsequently, became the practical standard for stable and efficient training across a wide range of tasks. Despite their widespread use, direct controlled comparisons between the two families under identical experimental conditions remain limited. This thesis presents three experiments comparing value-based and actor-critic algorithms on classic-control and Atari environments, using default hyperparameters throughout and evaluating results with the \texttt{rliable} statistical framework. The results confirm that actor-critic methods are the more consistent default at typical training budgets, while value-based methods remain competitive and outperform actor-critic methods on specific tasks given sufficient training time. A key methodological finding is that two standard aggregate metrics applied to the same data produce opposite conclusions about which family performs better, highlighting the insufficiency of single-number comparisons in reinforcement learning research.
}

\end{abstracts}
%\end{abstractlongs}


% ---------------------------------------------------------------------- 

